% GENERATED by scripts/build-lecture-slides.py; do not edit. Portable profile: XeTeX, theme and fonts from the core.
\makeatletter\IfFileExists{core/beamer/beamerthememx.sty}{\def\mxcorepath{core/}\def\input@path{{core/beamer/}{core/}}}{\def\mxcorepath{../core/}\def\input@path{{../core/beamer/}{../core/}}}\makeatother
\documentclass[aspectratio=169,10pt,t]{beamer}
\usetheme{mx}
\usepackage{graphicx,booktabs}
\usepackage[french]{babel}
\setlength{\parskip}{.45em}
\hypersetup{pdfauthor={},pdfsubject={Lecture slides},pdfkeywords={},pdftitle={Systèmes de numération}}
\providecommand{\N}{\mathbb{N}}
\providecommand{\Z}{\mathbb{Z}}
\providecommand{\D}{\mathbb{D}}
\providecommand{\Q}{\mathbb{Q}}
\providecommand{\R}{\mathbb{R}}
\providecommand{\C}{\mathbb{C}}
\providecommand{\vect}[1]{\overrightarrow{#1}}
\providecommand{\vi}{\vec{\imath}}
\providecommand{\vj}{\vec{\jmath}}
\providecommand{\vk}{\vec{k}}
\providecommand{\norm}[1]{\left\lVert #1\right\rVert}
\providecommand{\vc}[2]{\begin{pmatrix}#1\\#2\end{pmatrix}}
\providecommand{\deter}[2]{\det\left(#1,#2\right)}
\providecommand{\abs}[1]{\left\lvert #1\right\rvert}
\providecommand{\dd}{\,\mathrm{d}}
\providecommand{\e}{\mathrm{e}}
\providecommand{\Card}{\operatorname{Card}}
\providecommand{\E}{\operatorname{E}}
\providecommand{\V}{\operatorname{V}}
\providecommand{\degre}{^{\circ}}
\providecommand{\un}[1]{\,\mathrm{#1}}
\providecommand{\sci}[2]{#1\times10^{#2}}
\providecommand{\pgcd}{\operatorname{PGCD}}
\providecommand{\ppcm}{\operatorname{PPCM}}
\title{Systèmes de numération}
\subtitle{Cours et exemples}
\date{}
\author{}
\begin{document}
\mxtitleframe
\begin{frame}{Objectifs}
Lire une écriture en base entière et convertir un entier naturel d’une base à une autre.
\end{frame}
\begin{frame}{La valeur dépend de la position}
Dans une base entière \(b\ge2\), les chiffres sont \(0,1,\ldots,b-1\). L’écriture \((a_k\cdots a_1a_0)_b\) représente \(\sum_{i=0}^k a_i b^i\). Pour un entier non nul, le premier chiffre n’est pas nul; zéro s’écrit \((0)_b\). Pour les bases supérieures à dix, on peut convenir de lettres \(A=10\), \(B=11\), etc.
\end{frame}
\begin{frame}{La valeur dépend de la position}
\((1011)_2=1\times8+0\times4+1\times2+1=11\). \((203)_4=2\times16+0\times4+3=35\). Dans la base \(4\), le chiffre \(4\) est interdit.
\end{frame}
\begin{frame}{Convertir par divisions successives}
Diviser l’entier par \(b\), noter le reste, recommencer avec le quotient jusqu’à quotient nul. Lire les restes du dernier au premier.
\end{frame}
\begin{frame}{Convertir par divisions successives}
Pour \(45\) en base \(2\) :
\end{frame}
\begin{frame}{Convertir par divisions successives (2)}
\begin{center}\begin{tabular}{lll}
Nombre & Quotient & Reste \\ 
45 & 22 & 1 \\ 
22 & 11 & 0 \\ 
11 & 5 & 1 \\ 
5 & 2 & 1 \\ 
2 & 1 & 0 \\ 
1 & 0 & 1 \\ 
\end{tabular}\end{center}
\end{frame}
\begin{frame}{Convertir par divisions successives (3)}
Donc \(45=(101101)_2\). Vérification : \(32+8+4+1=45\).
\end{frame}
\begin{frame}{Convertir par divisions successives}
Chaque entier naturel possède une écriture unique en base \(b\ge2\) sans zéro initial, sauf l’écriture de zéro.
\end{frame}
\begin{frame}{Convertir par divisions successives}
La première division euclidienne impose le dernier chiffre, qui est l’unique reste dans \([0;b-1]\). Le quotient impose les suivants de la même manière. Pour un entier positif, les quotients successifs décroissent jusqu’à zéro. Les divisions fournissent donc une écriture finie et forcent tous ses chiffres.
\end{frame}
\begin{frame}{Changer de base et calculer}
En binaire, \((101)_2+(11)_2=(1000)_2\), car \(5+3=8\). Une retenue binaire correspond à deux unités du rang précédent, et non à dix.
\end{frame}
\begin{frame}{Changer de base et calculer}
Pour \(83\), \(83=5\times16+3\), \(16=5\times3+1\), \(3=5\times0+3\). Le premier reste donne les unités, le suivant les paquets de cinq, le dernier les paquets de vingt-cinq. Donc \(83=(313)_5\), vérifié par \(3\times25+1\times5+3=83\).
\end{frame}
\begin{frame}{Changer de base et calculer}
Aligner les unités, puis traiter chaque colonne de droite à gauche. Diviser la somme des chiffres et de la retenue par la base : le reste est le chiffre écrit, le quotient la retenue vers la colonne suivante. La retenue est un changement de rang, pas un chiffre supplémentaire ajouté au hasard.
\end{frame}
\begin{frame}{Changer de base et calculer}
Pour \((34)_5+(23)_5\), la colonne des unités donne \(4+3=7=5+2\) : écrire \(2\), retenir \(1\). La colonne suivante donne \(3+2+1=6=5+1\) : écrire \(1\) et retenir \(1\) devant. Le résultat est \((112)_5\). Contrôle décimal : \(19+13=32\) et \(25+5+2=32\).
\end{frame}
\begin{frame}{Changer de base et calculer}
Traduire d’abord l’écriture positionnelle en une équation portant sur \(b\). Résoudre cette équation, puis imposer \(b\in\mathbb{N}\), \(b\ge2\) et chaque chiffre strictement inférieur à \(b\). Une racine algébrique qui ne satisfait pas ces conditions n’est pas une base valide.
\end{frame}
\end{document}
